\documentclass[11pt,a4paper]{article}
% Preâmbulo comum. Não duplicar pacotes nos módulos.
\usepackage[utf8]{inputenc}
\usepackage[T1]{fontenc}
\usepackage[brazil]{babel}
\usepackage{amsmath,amssymb,amsthm,mathtools}
\usepackage{geometry}
\geometry{a4paper,margin=2.3cm}
\usepackage{microtype}
\usepackage{enumitem}
\usepackage{graphicx}
\usepackage{xcolor}
\usepackage{booktabs}
\usepackage{hyperref}
\hypersetup{colorlinks=true,linkcolor=blue!40!black,urlcolor=blue!50!black}

\setcounter{tocdepth}{2}
\numberwithin{equation}{section}

\theoremstyle{plain}
\newtheorem{teorema}{Teorema}[section]
\newtheorem{proposicao}[teorema]{Proposição}
\newtheorem{lema}[teorema]{Lema}
\newtheorem{corolario}[teorema]{Corolário}

\theoremstyle{definition}
\newtheorem{definicao}[teorema]{Definição}
\newtheorem{exemplo}[teorema]{Exemplo}
\newtheorem{exercicio}{Exercício}[section]

\theoremstyle{remark}
\newtheorem{observacao}[teorema]{Observação}
\newtheorem{contraexemplo}[teorema]{Contraexemplo}
\newtheorem{intuicao}[teorema]{Intuição}

\newcommand{\R}{\mathbb{R}}
\newcommand{\N}{\mathbb{N}}
\newcommand{\Z}{\mathbb{Z}}
\newcommand{\Q}{\mathbb{Q}}
\DeclareRobustCommand{\eps}{\varepsilon}
\DeclareRobustCommand{\vet}[1]{\vec{#1}}
\DeclareRobustCommand{\dif}{\mathrm{d}}
\DeclareRobustCommand{\abs}[1]{\left\lvert #1 \right\rvert}
\DeclareRobustCommand{\norm}[1]{\left\lVert #1 \right\rVert}
\DeclareMathOperator{\grad}{grad}
\DeclareMathOperator{\Div}{div}
\DeclareMathOperator{\rot}{rot}

\newcommand{\base}[1]{\par\medskip\noindent\textbf{Base para relembrar.} #1\par\medskip}
\newcommand{\cuidado}[1]{\par\medskip\noindent\textbf{Quando a hipótese falha.} #1\par\medskip}

\newcommand{\figura}[2]{%
  \begin{figure}[htbp]
  \centering
  \IfFileExists{#1}{%
    \includegraphics[width=0.84\textwidth]{#1}%
  }{%
    \fbox{\parbox{0.8\textwidth}{\centering Figura ainda não gerada:\\ \texttt{#1}}}%
  }
  \caption{#2}
  \end{figure}%
}

\newenvironment{esqueleto}{%
  \par\medskip\noindent\textbf{Esqueleto da demonstração.}\par\nopagebreak
  \begin{enumerate}[leftmargin=1.4em]
}{\end{enumerate}\medskip}

\title{Gabarito --- Derivadas}
\author{}
\date{}
\begin{document}
\maketitle
\tableofcontents

A ordem \'e a da se\c{c}\~ao de exerc\'icios de \texttt{03\_derivadas.tex}.

\section{B\'asico --- Pela defini\c{c}\~ao}

\textbf{Enunciado.} Calcule $(x^2)'$ em $x=1$ pelo quociente de Newton.

\begin{proof}
Para $h\neq 0$,
\[
\frac{(1+h)^2-1^2}{h}=\frac{1+2h+h^2-1}{h}=\frac{2h+h^2}{h}=2+h.
\]
Quando $h\to 0$, $2+h\to 2$. Logo $f'(1)=2$. Dado $\eps>0$, $\delta=\eps$ testemunha $\abs{(2+h)-2}=\abs{h}<\eps$.
\end{proof}

\section{B\'asico --- Tangente}

\textbf{Enunciado.} Reta tangente a $y=x^2$ em $(1,1)$.

\begin{proof}
Pelo exerc\'icio anterior, a inclina\c{c}\~ao \'e $2$ e o ponto de tang\^encia tem ordenada $1$. A reta \'e
\[
y=1+2(x-1),
\]
isto \'e, $y=2x-1$.
\end{proof}

\section{B\'asico --- Constante}

\textbf{Enunciado.} A derivada de $f(x)=7$ \'e $0$.

\begin{proof}
Para $h\neq 0$ e qualquer $a$, $(f(a+h)-f(a))/h=(7-7)/h=0$. O limite da fun\c{c}\~ao constante $0$ \'e $0$. Logo $f'(a)=0$.
\end{proof}

\section{B\'asico --- Linear}

\textbf{Enunciado.} A derivada de $f(x)=5x+2$ \'e $5$.

\begin{proof}
Para $h\neq 0$,
\[
\frac{f(a+h)-f(a)}{h}
=\frac{\bigl(5(a+h)+2\bigr)-(5a+2)}{h}
=\frac{5h}{h}=5.
\]
O quociente \'e a constante $5$, cujo limite \'e $5$. Portanto $f'(a)=5$ em todo $a$.
\end{proof}

\section{Intermedi\'ario --- Produto}

\textbf{Enunciado.} Calcule $(x^2\sin x)'$.

\begin{proof}
Sejam $u(x)=x^2$ e $v(x)=\sin x$. Sabemos $u'(x)=2x$ e $v'(x)=\cos x$. A regra do produto d\'a
\[
(uv)'(x)=u'(x)v(x)+u(x)v'(x)=2x\sin x+x^2\cos x.
\]
\end{proof}

\section{Intermedi\'ario --- Cadeia}

\textbf{Enunciado.} Calcule a derivada de $\cos(3x+1)$.

\begin{proof}
A fun\c{c}\~ao de fora \'e $\cos$, cuja derivada \'e $-\sin$. A fun\c{c}\~ao de dentro \'e $g(x)=3x+1$, com $g'(x)=3$. A regra da cadeia d\'a
\[
\frac{\dif}{\dif x}\cos(3x+1)=-\sin(3x+1)\cdot 3=-3\sin(3x+1).
\]
\end{proof}

\section{Intermedi\'ario --- Escada}

\textbf{Enunciado.} Escada de $5\,\mathrm{m}$, base a $0{,}3\,\mathrm{m/s}$, instante $x=4\,\mathrm{m}$. Calcule $y'$ e o instante em que $-x/y$ degenera.

\begin{proof}
A restri\c{c}\~ao \'e $x^2+y^2=25$. Com $y\neq 0$,
\[
y'=-\frac{x}{y}\,x'.
\]
No instante $x=4$ e $x'=3/10$, tem-se $y^2=25-16=9$. A ponta est\'a acima do ch\~ao, logo $y=3$, n\~ao $-3$. Ent\~ao
\[
y'=-\frac{4}{3}\cdot\frac{3}{10}=-\frac{4}{10}=-\frac{2}{5}.
\]
A unidade \'e $\mathrm{m/s}$: a ponta desce a $2/5\,\mathrm{m/s}$.

A express\~ao $-x/y$ deixa de fazer sentido quando $y=0$. A\'i $x^2=25$; do lado positivo do muro, $x=5$. A equa\c{c}\~ao derivada $2x x'+2y y'=0$ vira $10 x'=0$, logo exige $x'=0$. Uma base que ainda se move n\~ao pode permanecer na circunfer\^encia nesse instante.
\end{proof}

\section{Avan\c{c}ado --- Continuidade}

\textbf{Enunciado.} Deriv\'avel em $a$ implica cont\'inua em $a$.

\begin{proof}
Seja $f'(a)=\lim_{h\to 0}(f(a+h)-f(a))/h$. Para $h\neq 0$,
\[
f(a+h)-f(a)=\left(\frac{f(a+h)-f(a)}{h}\right) h.
\]
O primeiro fator tende a $f'(a)$ e o segundo tende a $0$. O produto dos limites \'e $f'(a)\cdot 0=0$. Logo $\lim_{h\to 0}f(a+h)=f(a)$, que \'e a continuidade de $f$ em $a$.
\end{proof}

\section{Avan\c{c}ado --- Cont\'inua e n\~ao deriv\'avel}

\textbf{Enunciado.} $f(0)=0$, $f(x)=x\sin(1/x)$ para $x\neq 0$: cont\'inua em $0$, sem derivada em $0$.

\begin{proof}
De $\sin^2 t+\cos^2 t=1$ segue $\abs{\sin t}\le 1$. Para $x\neq 0$,
\[
\abs{f(x)-f(0)}=\abs{x\sin(1/x)}\le \abs{x}.
\]
Dado $\eps>0$, escolha $\delta=\eps$. Se $\abs{x}<\delta$, ent\~ao $\abs{f(x)-f(0)}<\eps$. Isso \'e a continuidade em $0$.

Para $h\neq 0$,
\[
\frac{f(h)-f(0)}{h}=\frac{h\sin(1/h)}{h}=\sin\frac{1}{h}.
\]
Sejam, para cada inteiro $n\ge 1$,
\[
h_n=\frac{1}{\pi/2+2\pi n}, \qquad k_n=\frac{1}{3\pi/2+2\pi n}.
\]
Ent\~ao o quociente vale $\sin(\pi/2)=1$ em $h_n$ e $\sin(3\pi/2)=-1$ em $k_n$, e $0<h_n<1/(2\pi n)$, $0<k_n<1/(2\pi n)$. Suponha que o limite do quociente seja $L$. Para $\eps=1$ existe $\delta>0$ tal que $0<\abs{h}<\delta$ implica $\abs{\sin(1/h)-L}<1$. Escolha um inteiro $n>1/(2\pi\delta)$, o que existe porque os inteiros positivos s\~ao ilimitados. Ent\~ao $h_n$ e $k_n$ est\~ao em $(0,\delta)$, logo $\abs{1-L}<1$ e $\abs{-1-L}<1$. Portanto
\[
2=\abs{1-(-1)}\le \abs{1-L}+\abs{-1-L}<2.
\]
O quociente n\~ao tem limite, e $f'(0)$ n\~ao existe.
\end{proof}

\end{document}
